\scanpage{232}
\mspage{182}
% L1–2
et sûrement \deleted{(en dehors de toute hypothèse du type $b_n$ ou autre)} l'application canonique \inserted{\unsure{directe|dérivée|déduite} ``de localisation sur $Y$'', suivie de\ldots}
% L3
\[
H^n(X,F_X)\to H^0(Y,H^n(Z,F_Z))\to\mathrm{Hom}_\pi(\pi_n(Z),M)
\tag*{($\alpha$)}
\]
% L4–9
n'est autre que la \deleted{\unsure{$\simeq\mathrm{Hom}_{\pi}($}} composée des applications canoniques ci-dessus. Dans le cas actuel où on suppose $b_n$, la deuxième flèche dans ($\alpha$) est iso, donc la composée ``n'est autre'' que la flèche de localisation elle-même.
% L9–14
D'autre part $\pi_n(Z,z)\to\pi_n(X,x)$ dans ($*$) est nul, donc aussi $\pi_n(Z,\unsure{z|?})\to H_n(\unsure{X|Y},\mathbb Z)$, donc aussi la flèche composée ($*'$) qui s'en déduit. Revenant alors à la suite exacte de Leray (page précédente), on trouve que $H^n(\unsure{Y|X},F)\xrightarrow{\sim}H^n(X,\deleted{\unsure{f^*}}F_X)$.

% L15–19
\deleted{\unsure{Re}} Comment récupérer le cas où $F$ n'est \uline{pas} constant ? Dans le cas $n=2$, cela marchait, \deleted{\unsure{qui}} en passant \unsure{aux} rev. universels $\widetilde X,\widetilde Y$ de $X,Y$ \inserted{(pour pts $x,y$ choisis)}, et en écrivant les suites spectrales
% L20–21
\diagram{\begin{tikzcd}[ampersand replacement=\&,column sep=small]
H^*(Y,F)\arrow[d] \& E^{pq}_2\deleted{\unsure{\simeq}}\!=\deleted{\unsure{H}}H^p(\pi,H^q(\widetilde Y,F_{\widetilde Y})) \arrow[l,Rightarrow]\arrow[d]\\
H^{\unsure{*|n}}(X,F_X) \& E^{\prime pq}_2=H^p(\pi,H^q(\widetilde X,F_{\widetilde X})) \arrow[l,Rightarrow]
\end{tikzcd}}
% L22–24
L'hypothèse $b_n$ est conservée en passant de $f:X\to Y$ à $\widetilde f=\widetilde f_x:\widetilde X\unsure{-|\to}\widetilde Y$, donc on sait que
% archive sheet number: 232
